\documentclass{article}
\usepackage[paperwidth=140cm,paperheight=90cm,margin=0cm]{geometry}
\usepackage{fontspec}
\setmainfont{TeX Gyre Heros}
\usepackage{graphicx,tikz,xcolor}
\usepackage[hidelinks]{hyperref}
\hypersetup{pdftitle={SuperSycophantic: Stress-Testing Frontier Models from Single- to Multi-Turn Sycophancy},pdfauthor={SuperSycophantic authors}}
\pagestyle{empty}
\setlength{\parindent}{0pt}
\definecolor{heading}{HTML}{772725}
\definecolor{rulecolor}{HTML}{70518A}
\definecolor{muted}{HTML}{4B555B}
\newcommand{\at}[5]{\node[anchor=north west,inner sep=0pt,text width=#3cm,font=\fontsize{#4}{1.18\dimexpr#4pt\relax}\selectfont] at (#1,-#2) {#5};}
\newcommand{\plotat}[5]{\node[anchor=north west,inner sep=0pt] at (#1,-#2) {\includegraphics[width=#3cm,height=#4cm,keepaspectratio]{#5}};}
\newcommand{\cropat}[8]{\begin{scope}[shift={(#1,-#2)}]\clip (0,0) rectangle (#3,-#4);\node[anchor=north west,inner sep=0pt] at (-#6,#7) {\includegraphics[width=#5cm]{#8}};\end{scope}}
\newcommand{\sectionat}[4]{\at{#1}{#2}{#3}{40}{\color{heading}\bfseries #4}}

\begin{document}
\null
\begin{tikzpicture}[remember picture,overlay,x=1cm,y=1cm,shift={(current page.north west)}]
  % Header: project identity and author information stay above the figures.
  \plotat{2.4}{1.3}{11.0}{11.0}{assets/supersycophantic-logo.png}
  \at{15}{1.7}{105}{88}{\bfseries SuperSycophantic}
  \at{15}{5.9}{105}{43}{\bfseries Stress-Testing Frontier Models from Single- to Multi-Turn Sycophancy}
  \at{15}{9.2}{105}{26}{
    Terry Jingchen Zhang\textsuperscript{1,3,4,5}, Oscar Yasunaga\textsuperscript{2}, Wenyuan Jiang\textsuperscript{5}, Jessica Bo\textsuperscript{2,3}, Florent Draye\textsuperscript{9}, Aydin Javadov\textsuperscript{5},\\[5pt]
    Arth Singh\textsuperscript{2,6,7}, Lucille Muir\textsuperscript{2}, Victoria Oldemburgo De Mello\textsuperscript{2}, Bernhard Schölkopf\textsuperscript{5,8,9}, Zhijing Jin\textsuperscript{2,3,4,9}}
  \at{15}{13.0}{105}{22}{\color{muted}
    \textsuperscript{1}University of Oxford\quad \textsuperscript{2}University of Toronto\quad \textsuperscript{3}Vector Institute\quad \textsuperscript{4}EuroSafeAI\quad \textsuperscript{5}ETH Zürich\\[4pt]
    \textsuperscript{6}AIM Intelligence\quad \textsuperscript{7}Georgia Institute of Technology\quad \textsuperscript{8}ELLIS Institute Tübingen\quad \textsuperscript{9}Max Planck Institute for Intelligent Systems, Tübingen}
  \plotat{123}{1.8}{14.4}{6.5}{assets/neurips-logo.png}
  \at{124.2}{8.5}{13}{29}{\bfseries 2026 · Sydney}
  \node[anchor=north west,inner sep=0pt] at (126.5,-10.1) {\href{https://github.com/AI4Collaboration/SuperSycophantic}{\includegraphics[width=6.5cm]{assets/code-qr.pdf}}};
  \at{126.6}{16.6}{10}{22}{Code + data}
  \draw[rulecolor,line width=2.3pt] (2.5,-18.5) -- (137.5,-18.5);
  \draw[rulecolor!45,line width=1.1pt] (47,-20) -- (47,-82.5);
  \draw[rulecolor!45,line width=1.1pt] (93,-20) -- (93,-82.5);

  % Row one: benchmark, first-turn results, model-specific response to pressure.
  \sectionat{2.5}{20.0}{42.5}{One benchmark, two conversation regimes}
  \at{2.5}{23.8}{42.5}{29}{\bfseries 200 OBJ questions · 100 SUB decisions · 9 models}
  \cropat{2.5}{26.2}{21.8}{16.85}{43}{0.5}{0}{assets/benchmark-overview.png}
  \cropat{25}{26.2}{20}{10.15}{43}{22.55}{0}{assets/benchmark-overview.png}
  \at{25.2}{37.8}{19.3}{25}{\bfseries Escalating pressure}
  \at{25.2}{39.5}{19.3}{26}{Mild → Moderate → Strong}
  \at{25.2}{42.1}{19.3}{25}{\bfseries Changing influence family}
  \at{25.2}{43.8}{19.3}{26}{Family 1 → Family 2 → Family 3}
  \at{2.5}{47.0}{42.5}{28}{User framing → Initial answer → Follow-up pressure}
  \at{2.5}{50.2}{42.5}{27}{\color{muted}Static: fixed prompts. Adaptive: responds to the model.}

  \sectionat{49}{20.0}{42}{First-turn framing shifts answers}
  \plotat{49}{24.0}{42}{25.0}{assets/context.pdf}
  \at{49}{50.2}{42}{27}{\color{muted}OBJ: factual accuracy. SUB: selection of the user's option.}

  \sectionat{95}{20.0}{42.5}{Stronger pressure can reduce switching}
  \plotat{95}{24.0}{42.5}{25.0}{assets/claude-tone.pdf}
  \at{95}{50.2}{42.5}{27}{\color{muted}Supplementary SUB cohort: one follow-up, seven families.}

  % Row two: visual outcome definitions and controlled intervention results.
  \sectionat{2.5}{54.0}{42.5}{Track answer changes across turns}
  \cropat{2.5}{58.8}{42.5}{15.6}{70}{27.45}{3.25}{assets/outcome-definitions.png}
  \fill[white] (2.5,-58.8) rectangle (5.5,-60.2);
  \at{2.5}{76.0}{42.5}{38}{\bfseries SUB: A → B → A \quad Switch, then return}
  \at{2.5}{80.8}{42.5}{27}{\color{muted}OBJ: correct to incorrect. SUB: a switch without new decision-relevant information.}

  \sectionat{49}{54.0}{42}{User certainty increases switching}
  \plotat{49}{58.0}{42}{21.5}{assets/certainty.pdf}
  \at{49}{80.8}{42}{27}{\color{muted}Fixed word count and punctuation. Matched triplets: 1,055 initially correct OBJ; 900 SUB.}

  \sectionat{95}{54.0}{42.5}{A recheck instruction halves switching}
  \plotat{95}{58.0}{42.5}{21.5}{assets/recheck.pdf}
  \at{95}{79.8}{42.5}{27}{\color{muted}Opus-4.6: 32 initially correct OBJ + 40 SUB items × 3 wordings. With an answer key, both branches correct 7/7 wrong answers.}

  \draw[rulecolor,line width=2.3pt] (2.5,-85.4) -- (137.5,-85.4);
  \at{2.5}{86.3}{135}{37}{\bfseries Stress-test multi-turn sycophancy before deployment, while preserving valid self-correction.}
\end{tikzpicture}
\end{document}
